\section{Research Methodology and Paradigm}
This thesis adopts a \emph{design-science research} (DSR) paradigm. In design science research, the primary aim is to create and evaluate innovative artefacts that address real problems, in contrast to purely theoretical or observational studies \citep{GregorHevner2013}. \citep{Hevner2004} describe the design-science paradigm as one that “seeks to extend the boundaries of human and organisational capabilities by creating new and innovative artefacts”. In practical terms, this means that knowledge is generated through the act of building and applying a solution: “knowledge and understanding of a problem domain and its solution are achieved in the building and application of the designed artifact”. The DSR approach thus emphasises both utility and knowledge contribution. It is well-suited to engineering a correlator system, since the research outcome is a concrete technological solution (the digital FX correlator) and its validation, rather than solely abstract theory. The correlator itself embodies design knowledge: it is a tailored instantiation of signal-processing algorithms and software-engineering methods for the specific context of the Nooitgedacht Radio Telescope. As such, it qualifies as a DSR artefact (a construct and its implementation) whose creation and evaluation will both solve a practical problem and produce research insight.

Design science research aligns with a pragmatic research philosophy, where the truth of knowledge is judged by its usefulness in context. In this case, we are confronted with a practical problem (the lack of a flexible software correlator for a four-antenna interferometer) and we seek to learn by doing. This is in keeping with Simon’s notion of the “sciences of the artificial” and with the DSR principle that solutions to \emph{wicked} real-world problems often require iterative design rather than one-shot experimentation. In fact, DSR is commonly applied in engineering and computing disciplines, where systems, algorithms or instruments are designed to meet specific needs \citep{GregorHevner2013}. By adopting DSR, this research explicitly frames the correlator as a technologically relevant artefact and treats the design and implementation process as an essential part of knowledge generation.

The specific methodology follows a structured DSR process model. We employ the six-activity Design Science Research Methodology (DSRM) of \citep{Peffers2007}. These activities are: (1) Problem Identification and Motivation, (2) Objectives of a Solution, (3) Design and Development of the artefact, (4) Demonstration, (5) Evaluation, and (6) Communication. Figure~\ref{fig:dsr_cycle} illustrates this process and its iterative flow. In the context of this thesis, the six activities are instantiated as follows:

\begin{enumerate}
     \item \textbf{Problem identification and motivation.} The research begins by defining the specific problem environment and explaining why a new correlator is needed. For instance, Nooitgedacht Radio Telescope lacks today a modern digital correlator; current hardware solutions might be inflexible or obsolete, limiting the telescope's scientific capabilities. This procedure involves exploring the field (the telescope specification, observation goals, and existing infrastructure) in order to determine challenges that the software correlator will have to overcome. It also involves engaging with stakeholders (engineers, astronomers) to justify the research. This exercise is done to validate that the problem is significant and solving it will be worth it. For DSR, an explicitly stated problem motivates the design effort and gets unequivocal on the target for the artifact \citep{Hevner2004}.
 
     \item \textbf{Objectives of a solution.} From the problem definition, we derive clear objectives and requirements for the correlator artefact. These objectives are both functional and performance-oriented. For example, the correlator should process four input antenna signals in real-time, handle a specified bandwidth and number of frequency channels, and produce correct complex visibilities for all six baselines of the array. It should also integrate with the telescope’s data acquisition system and meet practical constraints (e.g.\ run on available computing hardware). The objectives can be quantitative (e.g.\ throughput $\geq$ X samples/s, latency $\leq$ Y ms) or qualitative (e.g.\ modular, extensible design). They are informed by both the problem domain and theoretical considerations (such as the computational complexity of the FX algorithm). In the DSR methodology, these objectives form the benchmarks for success; they are derived “from the problem definition and knowledge of what is possible and feasible” \citep{Peffers2007} and will guide evaluation.
 
     \item \textbf{Design and development.} With the goals informing, the artefact is developed. This is where, in practice, the FX correlator software architecture is designed and then realised. The design phase may include selecting algorithms (e.g. FFT libraries, data parallelisation methods) and decomposing the problem (F-engine channeliser and X-engine multiplier) into modules. For example, we might employ a pipeline where data flow is managed by Python, calling optimised FFT routines (NumPy/FFTW) and cross-multiplication kernels (possibly via Numba or GPUs). The process of development involves coding, hardware configuration (GPUs, network I/O) as required, and iteration across the implementation. Throughout, rigorous software engineering is maintained: requirements are validated against objectives, code is tested incrementally, and design choices are recorded. The result is an operational realization of the digital correlator artefact. \citep{Peffers2007} note, in design science research the design activity produces artefacts that may take the form of models, methods or instantiations \citep{Peffers2007}; here the artefact is a concrete system built in Python (with supporting compiled components as needed).
 
     \item \textbf{Demonstration.} Once a prototype correlator is developed, we demonstrate its use in solving the problem at hand. This involves exercising the artefact in a realistic setting or simulation. For example, a demonstration might involve feeding synthetic multi-antenna data into the correlator and showing that the output visibilities match theoretical expectations (such as known fringe patterns for injected signals). Alternatively, one could record signals from the actual antennas into a buffer and process them offline to confirm end-to-end functionality. The demonstration phase may use experiments or simulations, but always aims to show that the artefact can address an instance of the problem. The chosen demonstration method should reflect the problem context; it might, for instance, include a short field trial with the telescope or a lab setup with signal generators. Successful demonstration is not proof of final quality, but it provides evidence that the artefact design is on the right track.
 
     \item \textbf{Evaluation.} In the evaluation step, the artefact’s performance and utility are rigorously assessed against the objectives set earlier. This is a critical step to measure how well the correlator meets its goals. Evaluation may involve quantitative experiments (e.g.\ measuring processing speed, resource usage, numerical accuracy) and qualitative assessment (e.g.\ user feedback on interface or ease of use). For instance, one could measure the correlator’s end-to-end latency and compare it to the target, or correlate recorded astronomical signals and compare the output to results from a trusted reference system. According to the DSRM, evaluation “observe[s] and measure[s] how well the artifact supports a solution to the problem” by “comparing the objectives of a solution to actual observed results” \citep{Peffers2007}. Thus we would compare the correlator’s actual throughput, accuracy, and stability to the benchmarks defined in the objectives. This may also involve stress-testing under high data rates or diverse conditions. The evaluation is both empirical (tests and measurements) and analytical (e.g.\ analysing algorithmic scalability or potential bottlenecks). The outcome of this step determines whether the artefact is effective. Crucially, if the correlator falls short of requirements, this step produces insights that lead to refinement. As described in the DSR literature, at the end of evaluation researchers may “iterate back to activity 3 to try to improve the effectiveness of the artifact” \citep{Peffers2007}. In practice, this thesis is likely to undergo several design-test-evaluate cycles until the correlator satisfies the performance objectives.
 
    \item \textbf{Communication.} The final activity is to communicate the research problem, the artefact and its results to relevant audiences. In this project, the primary communication is this thesis document, which describes the motivation, design, and evaluation of the correlator. It may also include publication of code, presentations at conferences, or journal articles. \citep{Peffers2007} emphasise that communication should cover “the problem and its importance, the artifact, its utility and novelty, [and] the rigor of its design”. Thus we will explicitly discuss how the correlator contributes novel design knowledge and meets the design-science criteria (utility, performance). Documentation and dissemination are crucial DSR activities: they allow other researchers and practitioners to understand and potentially reuse the design. They also constitute the scholarly output of the research.
\end{enumerate}

These steps constitute a nominal workflow, but in reality the process is iterative rather than strictly linear. Figure~\ref{fig:dsr_cycle} conceptually summarises the DSR cycle for this project. As shown, each design-and-build cycle is followed by demonstration and evaluation, and insights from evaluation often feed back into further design. For example, if testing reveals that the correlator’s timing jitter is too high, the design may be revised accordingly. \citep{Peffers2007} note that in DSR the researcher can “iterate back to activity 3” (design) after evaluation \citep{Peffers2007}. Indeed, it is expected that the correlator design will be refined through multiple iterations of coding and testing. Moreover, the DSRM recognises that research need not start at activity 1 in sequence; sometimes a demonstrator may exist first and the problem is retroactively analysed, or objectives may be set by an industry requirement \citep{GregorHevner2013}. In our case the process is problem-driven, but the flexibility of DSR means we remain open to starting from different angles if needed. The build–evaluate loop (Figure~\ref{fig:dsr_cycle}) ensures that the final artifact is both functional and well-understood.

\begin{figure}[ht]
  \centering
  \includegraphics[width=0.8\textwidth]{sections/Images/dsr_cycle.png}
  \caption{Conceptual model of the Design Science Research (DSR) process, showing iterative activities from problem definition through design, demonstration, evaluation and communication (adapted from \citealt{Peffers2007}).}
  \label{fig:dsr_cycle}
\end{figure}

In summary, a DSR approach is appropriate for this correlator project because it explicitly links the engineering practice of building a system with the creation of new knowledge. By following a recognised DSR methodology \citep{Hevner2004,Peffers2007}, the research ensures that the artefact is developed systematically (with documented rationale and theoretical grounding) and evaluated rigorously. This methodology also emphasises relevance: the artifact is designed specifically to meet the needs of the Nooitgedacht telescope, and its success is measured by real-world criteria. Thus, the Digital FX Correlator thesis will not only result in a working instrument, but will also yield validated design insights, both of which are communicated to advance the field.